\chapter{债券：世界上最大、最无聊、最要命的市场}
\chaptermeta{10}

\begin{ledetext}
股票天天上头条。债券没人聊。可全世界所有资产的价格，起点都是从这儿抄的。
\end{ledetext}

\section{8000 块的借条}

李敏 38 岁，在南京一家医疗器械公司做出纳。

去年同事老邵装修差钱，找她借 8000。她没好意思提利息，让他写了张条。

条上四样东西，少一样这事就说不清：还多少（8000）、什么时候还（2027 年 3 月 1 日）、中间给不给利息（不给）、谁欠的（老邵）。

把这四样标准化、编上号、允许转手，它就叫债券。市场上管它们叫面值、期限、票息、发行人。

关键在"允许转手"。

李敏借出去的 8000 是死的，她得等到明年三月。可要是这张条能卖给隔壁工位的王姐，她随时能把钱拿回来，这 8000 就活了。

要能卖，条子就得长得一模一样。中国的债券面值统一定成 100 元一张，同一批发行的每一张都是同样的票息、同一个到期日。个性化的借条没法交易，标准化的借条才有市场。

全世界规模最大的市场，就是从这么一张纸长出来的。它每天成交以万亿美元计，买卖发生在机构之间的终端和电话里，没有涨跌停板，没有人在饭桌上聊它。

无聊，但要命。

\begin{statrow}{4}
  \statitem{353}{万亿美元}{全球债务总额，2026 年 3 月底（IIF）}
  \statitem{305}{\%}{相当于全球 GDP 的比例}
  \statitem{40}{万亿美元}{美国国债规模，2026 年 8 月突破}
  \statitem{30}{万亿美元+}{其中可流通的市场规模}
\end{statrow}

\section{李敏亏了 1130 块}

2021 年秋天，理财经理跟她说债基比定存强一点，稳。她挑了一只中长期纯债基金的 C 类份额，销售页面上标着"风险等级 R2"，放进去 4 万 7。

2022 年 11 月的一个早上，她在 3 号线上打开 App。

绿的。亏了 1130。

她的第一反应是：不是说好到期还本付息吗，这钱亏在哪儿了？

亏在她提前看了一眼价格。

\begin{egbox}{举个栗子}
一张面值 100 元、票息 3\%、还有 1 年到期的债券。到期那天它会付给你 103 元，本金加利息，白纸黑字，一个字都不会改。

现在市场变了，新发行的同类债券票息给到 4\%。你手里这张老债还想按 100 元出手？没人接。人家花 100 块买新债，一年后拿回 104；买你的老债，一年后只有 103。凭什么给你同样的价。

那老债该卖多少？倒着算。买家要的是 4\% 的回报，他愿意出的价格 P 得满足 P × 1.04 = 103，于是 P = 103 ÷ 1.04 = \textbf{99.04 元}。

票息没变，本金没变，发行人也没变。变的是别人那边的价码。

\end{egbox}

\textbf{价格和收益率是同一件事的两种说法，永远反向。}收益率上去，价格必然下来。这不是什么市场规律，这是一个除法。你没办法反对一个除法。

1 年期只跌了 9 毛 6，因为李敏只需要忍受一年的票息偏低。要是这张债还有 10 年到期，她得忍十年，补偿就得给够。

\begin{tablewrap}
\begin{xltabular}{\linewidth}{>{\raggedright\arraybackslash}p{0.20\linewidth}Xrr}
\toprule
\bthead{剩余期限} & \bthead{利率 3\% 时的价格} & \bthead{利率升到 4\% 后} & \bthead{跌幅} \\
\midrule
\textbf{1 年} & 100.00 元 & 99.04 元 & −0.96\% \\
\textbf{3 年} & 100.00 元 & 97.22 元 & −2.78\% \\
\textbf{5 年} & 100.00 元 & 95.55 元 & −4.45\% \\
\textbf{10 年} & 100.00 元 & 91.89 元 & −8.11\% \\
\textbf{30 年} & 100.00 元 & 82.71 元 & −17.29\% \\
\bottomrule
\end{xltabular}
\end{tablewrap}

同样一个百分点，30 年期的损失是 1 年期的十八倍。

\begin{pullquote}{}
债券本身没有变，是折算它的那把尺子变了。
\end{pullquote}

李敏买的那只基金持有的是一篮子中长期信用债。2022 年秋天资金面收紧，理财赎回一波接一波，价格一路往下走。

她没踩雷，也没碰上任何一笔违约。她只是买了一个每天都会被重新报价的东西，然后在最难看的那天早上打开了手机。

卖给她的时候没人讲过这一段。这不怪李敏，是话术的问题。

\section{硅谷银行是怎么死的}

\term{久期}{duration}\index{031@久期}可以粗糙地理解成：你平均要等多久，才能把钱收回来。

一张 10 年期的零息债，所有钱都在第 10 年到手，久期就是 10 年。一张每年付息的 10 年期债，中间陆续有钱回来，平均等待时间短一些，久期通常在 8 年出头。

它真正的用处是当尺子：\textbf{久期 8 年，市场利率每升 1 个百分点，价格大约跌 8\%。}久期 2 年就跌 2\%。粗暴，但误差不大，全世界的债券交易员每天靠这一条心算。

2023 年 3 月 10 日，硅谷银行被接管。从股价崩盘到关门，用了两天。

它买的债，本金一分钱都没损失。

\begin{egbox}{举个栗子：一家死于算术的银行}
2019 年底，硅谷银行的存款是 618 亿美元。两年后，1892 亿。它的客户是科技公司，那两年科技公司融资融到手软，钱全趴在它这儿。

钱来得太快，贷款放不出去那么多。它拿去买债：长期美国国债和抵押贷款支持证券，票息一个多点，久期普遍五年以上，光记在"持有至到期"科目下的就有 913 亿美元。

买的时候，联邦基金利率贴着零。到 2023 年初，美联储已经把它抬到 4.5\% 以上。

拿久期这把尺子量一下：4 个多百分点乘上五六年，账面浮亏就是二三十个百分点。2022 年底那本"持有至到期"的账，成本和市价之间差了 151 亿美元。

而硅谷银行全部的股东权益，163 亿。

会计规则允许这笔浮亏不进利润表，条件是你真能拿到到期。它拿不到。科技公司在烧钱，存款一天天往外走，银行必须卖债换现金。卖一笔，浮亏落地一笔。

市场看懂了这本账。48 小时里跑掉了几百亿存款。

\end{egbox}

\begin{warnbox}{小心}
银行给这件事起了个名字，叫期限错配。

把银行想成一个水池：进水管随时能被拧掉（存款当天就能取走），出水管却焊死在五年后（长期债券）。平时相安无事，因为不会所有人同一天来接水。这门生意的全部技术含量，就是赌"不会所有人同时来"。

这个比喻可以继续往下推。什么情况下大家会同时来？消息传得够快、取款够方便的时候。硅谷银行的储户是一个互相都认识的圈子，几个投资人在群里说一句"先把钱转出来"，手机银行三秒到账。两个条件同时满足了。

存款保险为什么能防挤兑？因为它把"我得赶在别人前面跑"这个念头从你脑子里拿掉。而硅谷银行的存款，超过九成在保险上限之外。

\end{warnbox}

\section{"无风险"的意思很窄}

国债收益率被叫做\term{无风险利率}{risk-free rate}\index{067@无风险利率}。

主权政府用本币借钱，几乎不会赖账。它有征税权，实在不行还能让央行下场接盘。所以它的\textbf{名义违约风险}极低。

"无风险"三个字的全部含义，到这儿就没了。

它不豁免通胀。你锁定 3\% 的票息，第二年通胀跑到 4.2\%，没有任何人赖你的账，但你每年悄悄穷 1.2 个点。

它更不豁免利率。李敏和硅谷银行手里拿的都是不会违约的债，一个亏掉 1130 块，一个亏掉了整家银行。

所以"债券保本"这句话是错的。不是不够严谨，是错的。

\begin{mythbox}{常见误解}
\mvfrow{误解}{"债券是保本的，买债基总不至于亏吧。"}
\mvfrow{事实}{2022 年是全球债券市场几十年来最惨的一年，美国长期国债那一年跌了 30\% 以上，比同期股票还难看。持有到期确实能拿回本金和票息，前提是发行人不违约。可\textbf{债券基金没有到期日}，它每天按市价估值、随时申赎。李敏买的不是一份合同，是一段价格。把"到期还本"和"随时能卖"混着讲，是这行的老毛病。}
\mvfrow{误解}{"这国国债收益率有 8\%，说明它信用好、给得起。"}
\mvfrow{事实}{恰恰可能相反。收益率是市场\emph{要求}的补偿，要得越狠，说明市场越不放心。2010 到 2012 年希腊 10 年期国债收益率一度冲到 30\% 以上，那不是希腊信用变好了。看一国的债得同时看两个数：收益率的绝对水平，和它相对于同期限最安全那条曲线的\textbf{利差}。}
\end{mythbox}

\section{打分的人，买单的人}

不是所有借条都由政府写。公司也发债，中国还有一类很特殊的城投债。

评级公司给每一张债打分。AAA 的意思接近"基本不会赖账"，BBB− 以上统称投资级，BB+ 往下叫高收益债，行里的叫法更实在：垃圾债。CCC 的含义大致是"这家公司现在就有麻烦"。

评级不是道德评价，是概率翻译。投资级债券的长期累计违约率在个位数，投机级能到两位数。你多要的那几个点利差，就是给这个概率开的价：\textbf{信用利差 = 你要的收益率 − 同期限国债收益率}。

\begin{databox}{一个十几年没改动的毛病}
给评级机构付钱的，是发债的公司，不是买债的人。你交学费，你考试，你还能挑老师。2008 年那批被打上 AAA、后来血本无归的次贷证券，就是这个模式的产物。

危机之后各国加了不少监管，付费结构基本没动。我的判断是它也动不了：让买方付费听起来很对，可你花钱买来的评级，转身就被市场上所有人免费看走了，没人愿意当那个掏钱的冤大头。这道题的名字叫搭便车，经济学解了一百年也没解开。

\end{databox}

城投债是中国债市的特产。地方政府过去通过融资平台公司举债搞基建，这些债法律上是企业债，市场却一直按"政府会不会兜底"来定价。这几年在做的事叫债务化解：把高息、短期、藏着的债，换成低息、长期、摆在明面上的债。利息负担确实降了，本金还在那儿，只是换了张脸和一个新的还款日。

这么多债，最后都堆在谁手里？

\begin{flowlist}
  \flowitem{01}{保险公司和养老金}{它们欠的是几十年后要付出去的赔款和养老金，久期长得吓人。要匹配就必须买长债。利率下行的时候它们最难受，因为负债的现值涨得比资产还快。}
  \flowitem{02}{银行}{拿存款去买债，赚一个利差，顺便满足流动性监管指标。硅谷银行干的就是这件事，只是干过了头。}
  \flowitem{03}{央行}{不为赚钱。它买卖国债是为了把政策意图传出去，价格是手段，不是目的。}
  \flowitem{04}{基金和 ETF}{替普通人打包，每天按市价估值。李敏买的是这一类，所以她能天天看见浮亏。}
  \flowitem{05}{外国投资者}{两个动机：储备要安全，利差要赚。这两个动机的相对分量正在变，下一节说。}
\end{flowlist}

这堆买家按行为分两类。保险、养老金、银行的自营账户是\textbf{配置盘}，买了就放着，只关心收益够不够覆盖负债成本。对冲基金、券商自营、一部分公募是\textbf{交易盘}，只关心三个月内价格往哪儿走。

债市剧烈波动的时候，动的一般不是配置盘，是交易盘在互相踩脚。配置盘的钱慢，交易盘的钱快，而价格永远由最快的那批钱定。

\section{第一把交椅换人了}

欧洲央行 2026 年 6 月的《欧元的国际角色》报告里，有一组数字。

\begin{barfig}
  \barrow{黄金（2024 年底）}{20}{20\%}
  \barrow{黄金（2025 年底）}{27}{27\%}
  \barrow{美国国债（2024 年底）}{25}{25\%}
  \barrow{美国国债（2025 年底）}{22}{22\%}
  \barrow{其他美元资产（2025 年底）}{20}{20\%}
  \barrow{欧元（2025 年底）}{15}{15\%}
\end{barfig}

截至 2025 年底，黄金占全球央行储备的 27\%，美国国债占 22\%。\textbf{黄金第一次超过美债，坐上了各国央行的第一把交椅。}一年之前，这个对比还是黄金 20\%、美债 25\%。

美元计价资产合计仍占 42\%，还是最大的单一币种阵营。所以这不是塌方，是配比在挪。为什么挪，第 12 章讲，那是一个关于"钱会不会被没收"的故事。

\begin{databox}{2026 数据}
IIF 在报告里警告，美国政府债务"日益走上不可持续的道路"。同一时期，外国投资者对美债的净买入总体保持稳定，但同步加配了日本和欧洲的主权债。这是分散化，不是撤离。

2025 年，欧元计价的国际债券发行量增长 30\%，接近 1 万亿欧元的历史新高。2026 年 8 月底，美国 10 年期国债收益率升到 2025 年 1 月以来的最高水平。

\end{databox}

\section{凭什么}

你手里任何一样东西的价格，本质上都是未来的现金流除以一个折现率。第 6 章推过这个式子。

那个折现率的第一块砖，就是同期限的国债收益率。它是所有定价的海拔基准。

\begin{chainflow}
\chainnode{国债收益率}\chainarrow \chainnode{信用债定价}\chainarrow \chainnode{股票估值的折现率}\chainarrow \chainnode{房贷与企业贷款定价}\chainarrow \chainnode{跨境资本流向}\chainarrow \chainnode{李敏的基金净值}
\end{chainflow}

10 年期美债收益率动 50 个基点，纳斯达克的估值模型要重算，上海写字楼的资本化率要重算，新兴市场的资金流向要重算。李敏家的房贷定价基准，也吊在同一条链子上。

\begin{talkbox}
  \talkline{新闻}{"截至 8 月底，美国 10 年期国债收益率升至 2025 年 1 月以来最高水平。"}
  \talkline{李敏}{"这跟我有什么关系。"}
  \talkline{翻译}{全世界给"未来的钱"打折时用的那个除数变大了 → 靠远期现金流撑估值的东西一起变便宜 → 已经发出去的长久期债券价格往下走 → 保险公司和养老金的账面被重估 → 借美元的新兴市场成本抬高 → 三个月后，你那只债基的净值曲线上会多出一个不太好看的弯。}
\end{talkbox}

\begin{pullquote}{}
债市打个喷嚏，全世界都得感冒。区别只在于，交易员当天就知道自己感冒了，李敏要等三个月才从基金净值上看出来。
\end{pullquote}

这本书不给任何投资建议。但有一件事可以说死：债券基金会亏钱，这是它的机制决定的，不是它出了差错。

\begin{takeaway}{本章一句话}
债券是一张标准化、可转手的借条。它的价格和收益率永远反向，而全世界所有资产的定价，都从国债收益率这块砖起步。

\begin{itemize}
  \item 价格与收益率反向是一个除法，不是市场规律。期限越长，同样的利率变动砸出来的坑越深：1 个百分点，30 年期跌 17.29\%。
  \item 久期是这个坑的深度尺。久期 8 年，利率升 1 个点，价格约跌 8\%。硅谷银行 151 亿美元的浮亏对 163 亿美元的股东权益，死因写在这把尺子上。
  \item "无风险"只指名义违约风险，不管通胀，不管利率。
  \item 截至 2025 年底，黄金以 27\% 超过美债的 22\%，成为全球央行第一大储备资产；美元计价资产合计仍占 42\%。
\end{itemize}

\end{takeaway}

\begin{bridgebox}
\textbf{下一章}债券是"先说好还你多少"。可要是我不想要那个说死的数字，我想跟着一家公司一起赚、一起赔呢？那就得先把公司切开。
\end{bridgebox}

